%% Reference template — the author's REAL Elsevier elsarticle submission.
%% latex.assemble() should produce output matching THIS structure (CiBM journal format):
%%   - \documentclass[preprint,12pt]{elsarticle}
%%   - \begin{frontmatter}: \title, \author[..]{}, \ead, \cortext, \affiliation[..]{organization=,city=,country=},
%%     \begin{abstract}..\end{abstract} (<=250 words), \begin{keyword} a \sep b \end{keyword}
%%   - body \section{}\label{} flow CONTINUOUSLY (NO \clearpage between sections)
%%   - tables and figures collected at the END, each preceded by \clearpage (wide ones in \begin{landscape})
%%   - \bibliographystyle{elsarticle-num} + numbered \cite{a,b}
\documentclass[preprint,12pt]{elsarticle}
\usepackage{xcolor,color}
\usepackage{soul}
\usepackage{amssymb}
\usepackage{amsmath}
\usepackage{booktabs}
\usepackage{pdflscape}
\usepackage{makecell}
\usepackage{kotex}
\usepackage{multirow}

\journal{Computers in Biology and Medicine}

\begin{document}
\begin{frontmatter}

\title{Efficient approximation of vessel wall stress under pulsatile hemodynamics using a snapshot-based coupling framework}

\author[label1,label2]{Se Hyeog Kim}
\author[label1]{Hyun Jin Kim\corref{cor1}}
\ead{kim.hyunjin@kaist.ac.kr}
\cortext[cor1]{Corresponding author}
\affiliation[label1]{organization={Department of Mechanical Engineering, Korea Advanced Institute of Science and Technology},
            city={Daejeon},
            country={South Korea}}

%% abstract (247/250 words) — confirms the ~250-word limit
\begin{abstract}
Fluid--structure interaction (FSI) analysis is essential for evaluating vessel-wall stress ... (<=250 words).
\end{abstract}

\begin{keyword}
Fluid--structure interaction \sep Vessel wall stress \sep Reduced-order modeling \sep Pulsatile hemodynamics \sep Snapshot-based coupling
\end{keyword}

\end{frontmatter}

%% Body sections flow continuously — NO \clearpage between them.
\section{Introduction}
\label{Introduction}
... \cite{a, b} ...

\section{Methods}
\label{Methods}
\subsection{Proposed snapshot-based coupling framework}
\subsubsection{One-dimensional reduced-order hemodynamics}
\begin{gather} ... \end{gather}

\section{Results}
\label{Results}
\subsection{Idealized LAD coronary artery model}

\section{Discussion}
\label{discussion}

\section{Conclusion}
\label{conclusion}

%% ===== END MATTER: tables then figures, each on its own \clearpage =====
\clearpage
\begin{table}[t]\centering
\begin{tabular}{@{}lr@{}}\toprule a & b \\ \bottomrule \end{tabular}
\caption{...}\label{tab:error}
\end{table}

\clearpage
\begin{figure}\centering
\includegraphics{Fig1.pdf}
\caption{...}\label{fig:fig1}
\end{figure}

\clearpage
\bibliographystyle{elsarticle-num}
\bibliography{sn-bibliography.bib}
\end{document}
